\section{Balls from holomorphic vanishing}
\label{sec:ball}

The aim of this section is to turn the order of a holomorphic zero into
the capacity of an embedded ball. We first fix the differential-form
normalization, then smooth a logarithmic singularity and transport the
resulting radial ball by Moser's deformation method.

Identify $\C^n$ with $\R^{2n}$ by writing $z=u+ix$.  For a smooth real
function $\phi$, use the convention
\[
 \dc\phi(X)=-\frac14\,d\phi(iX),
 \qquad \ddc\phi=d(\dc\phi).
\]
In particular, the standard primitive and symplectic form are
\begin{equation}
 \lambda_0=\dc|z|^2
   =\frac12\sum_{j=1}^n(u_j\,dx_j-x_j\,du_j),
 \qquad
 \omega_0=\ddc|z|^2=\sum_{j=1}^n du_j\wedge dx_j.
 \label{eq:ball-normalization}
\end{equation}
For every real tangent vector $X$, one has
\[
 \ddc\phi(X,iX)
 =\frac14\left.\Delta_{\zeta}\phi(z+\zeta X)\right|_{\zeta=0},
 \qquad \zeta\in\C.
\]
Thus a smooth function $\phi$ is plurisubharmonic precisely when these
quantities are nonnegative.  In that case $\ddc(|z|^2+\phi)$ is symplectic:
it is closed and its value on $(X,iX)$ is at least $|X|^2$.

We will also use that a smooth convex, coordinatewise nondecreasing
function of plurisubharmonic functions is plurisubharmonic.  Indeed, on
each complex line the ordinary chain rule gives
\[
 \Delta\Gamma(h_1,\ldots,h_N)
 =\sum_{j=1}^N\Gamma_j\Delta h_j
   +\sum_{i,j=1}^N\Gamma_{ij}
        \ip{\nabla h_i}{\nabla h_j}\ge0.
\]
The first sum is nonnegative by monotonicity, and the second is
nonnegative because the Hessian of $\Gamma$ is positive semidefinite.

\begin{theorem}[A ball from a holomorphic zero]
\label{thm:holomorphic-ball}
Let $n,m\ge1$, let $\Omega\subset\C^n$ be an open set containing $0$,
and let $f=(f_1,\ldots,f_m):\Omega\to\C^m$ be holomorphic.  Suppose that,
for some integer $k\ge1$,
\begin{enumerate}[label=\textup{(\roman*)}]
 \item $f^{-1}(0)=\{0\}$ and $|f(z)|=O(|z|^k)$ as $z\to0$;
 \item for every $0<s<1$, the set
 $\{z\in\Omega:|f(z)|^2\le s\}$ is a compact subset of $\Omega$.
\end{enumerate}
Set
\[
 \tau=|f|^2,\qquad V=\{z\in\Omega:\tau(z)<1\},
 \qquad \omega=\ddc(|z|^2+\tau).
\]
For every $0<c<\pi k$, there is a smooth symplectic embedding
\[
 \bigl(B^{2n}(c),\omega_0\bigr)\longrightarrow(V,\omega).
\]
\end{theorem}

\begin{proof}
\noindent\textit{Step 1: a logarithmic potential with the required slope.}
For a holomorphic tuple, differentiation along a complex line gives
\begin{align}
 \ddc\tau(X,iX)&=|df(X)|^2,\notag\\
 \ddc\log\tau(X,iX)
 &=\frac{|f|^2|df(X)|^2
       -\left|\sum_{j=1}^m\overline{f_j}\,df_j(X)\right|^2}{|f|^4}
 \quad\text{on }\Omega\setminus\{0\}.
 \label{eq:ball-log-levi}
\end{align}
The second expression is nonnegative by Cauchy--Schwarz.  Hence $\tau$
and, away from its zero, $\log\tau$ are plurisubharmonic.

Fix $0<a<k$; we will embed the ball of capacity $\pi a$.  Choose
\[
 \frac ak<b<1,\qquad a<\mu<bk.
\]
Choose numbers $\log b<t_-<t_+<0$ and a smooth nondecreasing cutoff
$\beta:\R\to[0,1]$ which is $0$ for $t\le t_-$ and $1$ for $t\ge t_+$.
Define $\chi$ on $(-\infty,0)$ by
\[
 \chi'(t)=b+\beta(t)(e^t-b),
\]
choosing its additive constant so that $\chi(t)=e^t$ for $t\ge t_+$.
This is possible since $\chi'=e^t$ there.  Moreover, $\chi'>0$ and
\[
 \chi''(t)=\beta'(t)(e^t-b)+\beta(t)e^t\ge0,
\]
because the transition occurs where $e^t>b$.  For $t\le t_-$, the
function $\chi$ is affine with slope $b$.

It follows that
\[
 H=\chi(\log\tau)
\]
is smooth and plurisubharmonic on $V\setminus\{0\}$ and satisfies
$H=\tau$ wherever $\tau\ge s_0:=e^{t_+}$.  The vanishing assumption gives
$\tau(z)\le C_0|z|^{2k}$ near zero.  Since $\chi$ is eventually affine,
there is a constant $C_1$ such that
\begin{equation}
 H(z)\le bk\log|z|^2+C_1
 \quad\text{for all sufficiently small }z\ne0.
 \label{eq:ball-singularity}
\end{equation}

\medskip
\noindent\textit{Step 2: smooth the singularity without changing the form near the boundary.}
The strict inequality $\mu<bk$ lets a smooth radial logarithm dominate
$H$ close to zero while remaining below it on an outer sphere.
Write $B_r=\{z\in\C^n:|z|<r\}$, and choose $r_2>0$ with
$\overline{B}_{r_2}\subset V$.  For $0<\delta\le1$, consider
\[
 P_\delta(z)=\mu\log(|z|^2+\delta)-M.
\]
For each fixed $r>0$, the radial calculation in Step~4 identifies
$(B_r,\ddc(|z|^2+P_\delta))$ with a standard ball whose squared radius
is $r^2+\mu r^2/(r^2+\delta)$. Since $\mu>a$, this exceeds $a$ for
all sufficiently small $\delta$. We will therefore arrange that the
modified potential agrees with $P_\delta$ on a ball whose radius is
fixed before $\delta$ is chosen.

The functions $P_\delta$ are plurisubharmonic by
Equation~\eqref{eq:ball-log-levi}, applied to the nowhere-zero
holomorphic tuple $(\sqrt\delta,z_1,\ldots,z_n)$.  Choose $M$ so that
\[
 M>\mu\log(r_2^2+1)-\min_{|z|=r_2}H(z)+2.
\]
The minimum is finite because $f$ has no zero on this sphere.  Then
$P_\delta<H-2$ on $|z|=r_2$ for every $0<\delta\le1$.
On the other hand, Equation~\eqref{eq:ball-singularity} gives
\[
 P_\delta(z)-H(z)
 \ge(\mu-bk)\log|z|^2-M-C_1.
\]
As $z\to0$, this lower bound tends to $+\infty$, uniformly in $\delta$.
We may therefore fix $0<r_1<r_2$, independently of $\delta$, such that
\begin{equation}
 P_\delta>H+2\quad\text{on }0<|z|\le r_1
 \quad\text{for every }0<\delta\le1.
 \label{eq:ball-inner-domination}
\end{equation}
Fix from now on one positive $\delta$ satisfying
\begin{equation}
 0<\delta<\min\left\{1,\ r_1^2\left(\frac\mu a-1\right)\right\}.
 \label{eq:ball-delta-choice}
\end{equation}
The upper bound is positive because $\mu>a$.

We next replace the singularity of $H$ by $P_\delta$ using a
regularized maximum, the standard gluing device in Richberg's smoothing
method; see \cite[Section~2]{HarveyLawsonPlis2020}. A related replacement
of K\"ahler potentials appears in
\cite[Proposition~3.3]{WittNystrom2018} in a projective setting.
We give the elementary two-function construction explicitly. Let
$\vartheta\in C_c^\infty((-1,1))$ be nonnegative and even, with
$\int_{\R}\vartheta=1$, and set
\[
 \rho(s)=\int_{\R}|s-t|\vartheta(t)\,dt,
 \qquad
 \operatorname{rmax}(h,p)=\frac{h+p+\rho(h-p)}2.
\]
The function $\rho$ is smooth and convex, with $|\rho'|\le1$, and
equals $|s|$ for $|s|\ge1$.  Thus $\operatorname{rmax}$ is smooth,
convex, nondecreasing in both arguments, and equals $\max(h,p)$ when
$|h-p|\ge1$.  Define $\widetilde H$ to be
$\operatorname{rmax}(H,P_\delta)$ on $B_{r_2}\setminus\{0\}$, to be
$H$ outside $B_{r_2}$, and to take the value $P_\delta(0)$ at zero.
The outer comparison gives equality with $H$ in a neighborhood of
$\partial B_{r_2}$, and Equation~\eqref{eq:ball-inner-domination} gives
equality with $P_\delta$ on $B_{r_1}$.  Consequently $\widetilde H$ is
smooth and plurisubharmonic on all of $V$.

The difference $\widetilde H-\tau$ has compact support in $V$.
Indeed, set $s_2=\max_{\overline B_{r_2}}\tau<1$, choose
$\max(s_0,s_2)<\sigma<1$, and put
\begin{equation}
 C=\{z\in\Omega:\tau(z)\le\sigma\}.
 \label{eq:ball-support}
\end{equation}
This is a compact subset of $V$ by hypothesis.  Outside $C$ a point
lies outside $\overline B_{r_2}$ and has $H=\tau$, so that
$\widetilde H-\tau$ vanishes there.

\medskip
\noindent\textit{Step 3: transport the smoothed form.}
We now use Moser's deformation argument \cite[Section~4]{Moser1965}.
The support estimate above supplies a complete flow even though $V$
need not be compact. Consider the symplectic forms
\[
 \omega_s=\ddc\bigl(|z|^2+(1-s)\tau+s\widetilde H\bigr),
 \qquad 0\le s\le1,
\]
and write $\widetilde\omega=\omega_1$.  Each form is symplectic because
the two non-Euclidean potentials are plurisubharmonic.  Define the
smooth time-dependent vector field $Y_s$ by
\[
 \iota_{Y_s}\omega_s=-\gamma,
 \qquad \gamma=\dc(\widetilde H-\tau).
\]
Both $\gamma$ and every $Y_s$ are supported in the fixed compact set
$C$.  Extending $Y_s$ by zero outside $V$ gives a smooth vector field
on $\C^n$, with its coefficients and first derivatives bounded
uniformly for $s\in[0,1]$.  The ordinary differential equation for
this field has a flow $\varphi_s$ throughout this time interval.
The ambient flow fixes $\C^n\setminus V$ pointwise and hence restricts
to diffeomorphisms of $V$.  This argument applies even if $V$ is
disconnected and places no regularity assumption on its boundary.

Since $d\omega_s=0$, differentiation of pullbacks gives
\[
 \frac{d}{ds}\varphi_s^*\omega_s
 =\varphi_s^*\bigl(d\gamma+d\iota_{Y_s}\omega_s\bigr)=0.
\]
Therefore
\begin{equation}
 \varphi_1^*\widetilde\omega=\omega,
 \qquad
 \varphi_1^{-1}:(V,\widetilde\omega)\longrightarrow(V,\omega)
 \text{ is symplectic}.
 \label{eq:ball-moser}
\end{equation}

\medskip
\noindent\textit{Step 4: read off the ball in the radial model.}
The deformation has reduced the problem to an explicit radial form.
On $B_{r_1}$ the new potential, up to the additive constant $-M$, is
$\Psi(|z|^2)$, where
\[
 \Psi(t)=t+\mu\log(t+\delta).
\]
Define the radial map
\[
 R(z)=\sqrt{\Psi'(|z|^2)}\,z
     =\sqrt{1+\frac\mu{|z|^2+\delta}}\,z.
\]
This is the radial case of the symplectic-coordinate construction in
\cite[proof of Theorem~1.1, Equation~(23)]{LoiZuddas2008}.
In the pullback of $\lambda_0$ from
Equation~\eqref{eq:ball-normalization}, the terms differentiating the
radial factor cancel.  Hence
\[
 R^*\lambda_0=\Psi'(|z|^2)\lambda_0=\dc\Psi(|z|^2),
 \qquad R^*\omega_0=\widetilde\omega\big|_{B_{r_1}}.
\]
The squared radial coordinate is strictly increasing, because
\[
 \frac{d}{dt}\bigl(t\Psi'(t)\bigr)
 =\frac{d}{dt}\left(t+\frac{\mu t}{t+\delta}\right)
 =1+\frac{\mu\delta}{(t+\delta)^2}>0.
\]
At zero the map is smooth and has derivative
$DR(0)=\sqrt{1+\mu/\delta}\,I$, so its inverse is smooth there as
well.  Thus $R$ is a diffeomorphism from $B_{r_1}$ onto the standard
ball $B^{2n}(\pi Q_\delta)$, where
\[
 Q_\delta=r_1^2+\frac{\mu r_1^2}{r_1^2+\delta}>a
\]
by Equation~\eqref{eq:ball-delta-choice}.  If
$\iota:B_{r_1}\hookrightarrow V$ denotes inclusion, the map
\[
 \varphi_1^{-1}\circ\iota\circ
       R^{-1}\big|_{B^{2n}(\pi a)}:
 B^{2n}(\pi a)\longrightarrow V
\]
is a smooth embedding, and its pullback of $\omega$ is $\omega_0$ by
Equation~\eqref{eq:ball-moser} and the radial pullback identity.
Since $a$ was arbitrary in $(0,k)$, the theorem follows.
\end{proof}
